\documentclass[a4paper]{article}

% --- Typography: Times New Roman Throughout ---
\usepackage{fontspec}
\setmainfont{Times New Roman}[Ligatures=TeX]

% --- Page Geometry ---
\usepackage[top=1.25cm,bottom=1.25cm,left=1.50cm,right=1.50cm]{geometry}
\usepackage{amsmath,amssymb,amsfonts}
\usepackage{booktabs}
\usepackage{array}
\usepackage{microtype}
\usepackage{caption}
\usepackage{titlesec}
\usepackage{url}

% Academic styling
\captionsetup{font=small,labelfont=bf,skip=2.5pt}
\titlespacing*{\section}{0pt}{3.5pt plus 0.8pt minus 0.5pt}{1.5pt plus 0.5pt minus 0.3pt}
\titleformat{\section}{\fontsize{10}{11.5}\bfseries}{\thesection.}{0.4em}{}

\setlength{\parskip}{1.2pt plus 0.4pt minus 0.2pt}
\setlength{\parindent}{1.0em}
\pagestyle{empty}

\begin{document}
\fontsize{8.8}{10.8}\selectfont

% --- Title, Authors, and Institutional Affiliation ---
\begin{center}
  {\fontsize{12.5}{14.2}\bfseries ORBIT-INTEL 6G: Intelligent Resource Management and Predictive\\[1.2pt]
  Reliability Orchestration for Integrated Non-Terrestrial Networks\par}
  \vspace{3pt}
  {\fontsize{9.8}{11.5}\bfseries Slock Ahuja\quad and\quad Dr. Praveen Kumar Sharma\par}
  \vspace{1.2pt}
  {\fontsize{8.5}{10.2}\itshape Department of Electronics and Communication Engineering, MIT-ADT University, Pune, India\par}
  \vspace{1.2pt}
  {\fontsize{7.8}{9.5}\rmfamily Project Showcase: IEEE Communications Society Student Tech Leadership Conference (STLC) Tunisia 2026\par}
  {\fontsize{7.8}{9.5}\rmfamily Repository: \url{https://github.com/SlockAhuja/orbit-intel-6g} $\,\cdot\,$ Showcase: \url{https://slockahuja.github.io/orbit-intel-6g/}\par}
\end{center}

\vspace{-3pt}
\begin{abstract}
\fontsize{8.2}{10.1}\selectfont
\noindent Integrated Space-Air-Ground Non-Terrestrial Networks (SAGIN) form an indispensable pillar of Sixth-Generation (6G) wireless systems, delivering pervasive connectivity across maritime, aeronautical, and remote wilderness regions. However, high orbital velocity ($v \approx 7.56\text{ km/s}$), brief visibility windows ($<8\text{ min}$), Doppler shifts ($\pm 48\text{ kHz}$), and atmospheric rain fade ($>15\text{ dB}$) severely challenge Ultra-Reliable Low-Latency Communication (URLLC). This paper presents \textbf{ORBIT-INTEL 6G}, a two-tier intelligence framework addressing these bottlenecks. At Tier~1, an intra-tier Proximal Policy Optimization (PPO) reinforcement learning agent governs dynamic user association, orthogonal subcarrier allocation, and adaptive power control. At Tier~2, Confidence-Aware Predictive Reliability Orchestration (CAPRO) performs multi-horizon link forecasting ($H \in \{1, 3, 5, 10\}\text{ s}$) with distribution-free split-conformal calibration ($1-\alpha=95\%$) and dual aleatoric-epistemic uncertainty estimation to execute risk-adjusted proactive handovers. Reported experimental results from an extensive 400-trial Monte Carlo benchmark campaign demonstrate a packet delivery ratio of \textbf{99.41\%}, a 99th-percentile tail latency of \textbf{6.82 ms} (meeting sub-10ms URLLC requirements), an outage probability of \textbf{0.59\%}, and an \textbf{89.3\%} reduction in handover ping-pong flapping relative to conventional 3GPP reactive baselines.
\end{abstract}

\vspace{0.5pt}
\noindent\textbf{\fontsize{8.2}{10.1}\selectfont Keywords}---{\fontsize{8.2}{10.1}\selectfont 6G, Non-Terrestrial Networks, SAGIN, Resource Management, PPO, Predictive Intelligence, CAPRO.}

\section{Introduction}
Non-Terrestrial Networks (NTN) have emerged as a foundational architecture in 3GPP Releases 17 and 18 to bridge the digital divide by integrating satellite, high-altitude platform, and terrestrial infrastructure into unified 6G communications. The resulting Space-Air-Ground Integrated Network (SAGIN) offers unprecedented coverage resilience across oceanic transport routes, polar corridors, and disaster zones. However, delivering deterministic URLLC guarantees over non-terrestrial links is hindered by three core physical limitations: rapid Low Earth Orbit (LEO) satellite motion inducing transient topocentric visibility windows ($<8\text{ minutes}$); severe Doppler dynamics ($\pm 48\text{ kHz}$ at Ka-band $28\text{ GHz}$); and dynamic atmospheric attenuation (ITU-R P.618) during convective precipitation events.

Crucially, standard cellular handover mechanisms (such as 3GPP Event A3/A4 triggers) operate purely reactively, initiating target node search only after signal degradation or radio link failure (RLF) has already manifested. In satellite-terrestrial feeder topologies characterized by round-trip propagation delays of tens of milliseconds, reactive re-association induces packet dropping, retransmission storms, and severe handover ping-pong flapping. To overcome these constraints, 6G NTN requires a paradigm shift toward uncertainty-aware proactive orchestration.

\section{Proposed ORBIT-INTEL 6G Framework}
The ORBIT-INTEL 6G architecture models a five-tier heterogeneous SAGIN hierarchy: Tier~1 comprises LEO satellites ($500\text{--}600\text{ km}$ altitude) providing wide-area backhaul and dynamic spot beams; Tier~2 incorporates stratospheric High-Altitude Platform Stations (HAPS, $20\text{ km}$) acting as quasi-stationary aerial base stations free from tropospheric rain fade; Tier~3 deploys tactical Unmanned Aerial Vehicles (UAVs, $0.3\text{--}2\text{ km}$) for localized line-of-sight restoration in obstructed areas; and Tiers~4 and~5 encompass terrestrial gNodeBs ($30\text{ m}$) and mobile user terminals distributed across five representative geographic environments (Urban Hotspots, Rural Farmlands, Maritime Corridors, Mountain Observatories, and Post-Disaster Emergency Zones). The platform couples a standard-compliant orbital propagation engine with a two-tier intelligence hierarchy: an intra-tier physical resource optimizer and a cross-tier predictive reliability orchestrator.

\section{AI-Based Resource Management}
Physical-layer radio resource allocation within active satellite footprints is orchestrated by an autonomous deep reinforcement learning agent based on Proximal Policy Optimization (PPO). Interacting with a standardized Gymnasium non-terrestrial environment, the agent observes the normalized state $s_t = [h_t, \gamma_t, B_t, P_t, \text{Alt}_t] \in \mathbb{R}^5$, representing channel gain, SINR, available bandwidth, transmit power budget, and satellite altitude. The actor-critic network selects discrete actions governing: (i)~\textit{Terminal Association}, weighting topocentric elevation and dwell time with a hysteresis bonus ($+4.5$) to prevent premature handoffs; (ii)~\textit{Orthogonal Subcarrier Slicing}, assigning orthogonal frequency partitions ($16\text{--}32$ channels) across adjacent spot beams to eliminate co-channel interference; and (iii)~\textit{Water-Filling Power Control}, allocating transmit power ($P_c \le P_{\max}$) to mitigate localized path attenuation. The policy optimizes the scalar reward:
\begin{equation}
  R_t = \text{Throughput}_t - \text{Delay}_t - 10 \cdot \text{PacketLoss}_t
\end{equation}
Heavily penalizing packet losses, the PPO agent reliably maintains maximal spectral efficiency during stable orbital passes.

\section{Predictive Reliability Orchestration}
While physical-layer PPO manages intra-tier spectrum and power slicing, network-layer continuity across constellation handovers requires predictive intelligence. ORBIT-INTEL 6G introduces Confidence-Aware Predictive Reliability Orchestration (CAPRO), operating as an independent network-layer framework. Crucially, PPO and CAPRO perform distinct functions: PPO manages intra-tier physical resources, whereas CAPRO orchestrates cross-tier proactive handovers and routing across LEO, HAPS, and UAV relays. They are complementary components rather than a single joint model.

CAPRO executes three sequential stages: (i)~\textit{Multi-Horizon Neural Forecasting}, ingesting rolling telemetry time-series to predict link quality trajectories over lookahead horizons $H \in \{1, 3, 5, 10\}\text{ s}$; (ii)~\textit{Split-Conformal Calibration}, deriving distribution-free prediction intervals $[\hat{y}_{t+H}^{\text{lower}}, \hat{y}_{t+H}^{\text{upper}}]$ from non-conformity residuals with finite-sample marginal coverage $\mathbb{P}(Y_{t+H} \in \mathcal{C}(X_t)) \ge 95\%$, decomposing variance into aleatoric channel noise $\sigma_{\text{aleat}}$ and epistemic model uncertainty $\sigma_{\text{epist}}$; and (iii)~\textit{Risk-Adjusted Utility Maximization}, ranking candidate target nodes $c \in \{\text{LEO}_2, \text{HAPS}, \text{UAV}\}$ via:
\begin{equation}
  \mathcal{U}(c) = \hat{\gamma}_c(t+H) - \lambda \cdot (\sigma_{\text{aleat}, c}(t+H) + \sigma_{\text{epist}, c}(t+H)) - C_{\text{handover}}
\end{equation}
where $\lambda = 1.5$ balances risk aversion and $C_{\text{handover}}$ penalizes signaling overhead. When the active link utility degrades, CAPRO triggers a proactive handover $1.8\text{ s}$ before line-of-sight failure, enabling seamless transition without packet loss.

\clearpage % ==================== STRICT TRANSITION TO PAGE 2 ====================

\section{Simulation and Experimental Methodology}
The ORBIT-INTEL 6G evaluation framework implements physical aerospace kinematics and propagation models adhering to international telecommunications standards:
\begin{itemize}\setlength{\itemsep}{1.2pt}\setlength{\parskip}{0pt}
  \item \textbf{Kinematics \& Orbit Propagation}: LEO circular orbits at $500\text{--}600\text{ km}$ altitude with minimum elevation threshold $\theta \ge 20^\circ$. Slant range distance is evaluated geometrically as $d(t) = \sqrt{R_E^2 \sin^2\theta + 2 R_E h + h^2} - R_E \sin\theta$, where $R_E = 6371\text{ km}$ is Earth's mean radius and $h$ is orbital altitude.
  \item \textbf{Path Loss \& Doppler Dynamics}: Free Space Path Loss obeys 3GPP TR 38.811 at Ka-band ($f_c = 28\text{ GHz}$): $\text{FSPL} = 20\log_{10}(d) + 20\log_{10}(f_c) - 147.55\text{ dB}$. Doppler shifts $f_D(t) = f_c \cdot (v_{\text{rel}}(t) / c)$ span $\pm 48\text{ kHz}$ across typical LEO satellite passes.
  \item \textbf{Hydrometeor Fading}: Atmospheric rain attenuation follows ITU-R P.618-13, implementing specific attenuation $\gamma_R = k R^\alpha\text{ dB/km}$ across precipitation rates from clear sky ($0\text{ mm/h}$) to severe convective downpours ($100\text{ mm/h}$).
  \item \textbf{Benchmark Baseline Configurations}: Three established schemes are evaluated: (M1)~\textit{Reactive Baseline} (standard 3GPP Event A3/A4 triggers with hysteresis and time-to-trigger timers); (M2)~\textit{Threshold-Based} (static $10\text{ dB}$ SNR margin); and (M3)~\textit{Prediction-Only} (neural point forecast without conformal uncertainty bounds).
\end{itemize}

\section{Experimental Evaluation}
The proposed architecture was evaluated across an exhaustive \textbf{400-trial Monte Carlo} simulation campaign under stochastic rain fading, dynamic multi-user mobility across five topological environments, and orbital satellite transitions. The reported experimental results are summarized in Table~\ref{tab:benchmarks}.

\begin{table}[h!]
  \centering
  \fontsize{8.2}{10.2}\selectfont
  \caption{Reported Experimental Benchmark Results (400-Trial Monte Carlo Campaign)}
  \label{tab:benchmarks}
  \renewcommand{\arraystretch}{1.12}
  \begin{tabular}{lccccc}
    \toprule
    \textbf{Orchestration Methodology} & \textbf{PDR (\%)} & \textbf{99\% Tail Latency} & \textbf{Outage Rate} & \textbf{Ping-Pong / hr} & \textbf{Energy (J/Mbit)} \\
    \midrule
    M1: Reactive Baseline (3GPP A3/A4) & $91.20 \pm 1.45$ & $38.45\text{ ms}$ & $8.80\%$ & $18.4 \pm 2.1$ & $4.85 \pm 0.35$ \\
    M2: Threshold-Based Heuristic & $94.65 \pm 1.10$ & $26.12\text{ ms}$ & $5.20\%$ & $12.1 \pm 1.6$ & $3.92 \pm 0.28$ \\
    M3: Prediction-Only (Point-Forecast) & $97.15 \pm 0.85$ & $14.80\text{ ms}$ & $2.95\%$ & $6.8 \pm 1.1$ & $3.15 \pm 0.22$ \\
    \textbf{M4: CAPRO (Proposed Framework)} & $\mathbf{99.41 \pm 0.32}$ & $\mathbf{6.82\text{ ms}}$ & $\mathbf{0.59\%}$ & $\mathbf{1.9 \pm 0.4}$ & $\mathbf{2.10 \pm 0.16}$ \\
    \midrule
    \textit{Measured Relative Improvement} & \textbf{+8.21\% Gain} & \textbf{-82.26\% Reduction} & \textbf{-93.30\% Reduction} & \textbf{-89.30\% Reduction} & \textbf{-56.70\% Efficiency} \\
    \bottomrule
  \end{tabular}
  \vspace{2pt}
  \parbox{\linewidth}{\fontsize{7.6}{9.2}\selectfont\textit{Note}: Reported experimental results from the verified 400-trial Monte Carlo benchmark campaign.}
\end{table}

The reported experimental data confirms four pivotal performance findings:
\begin{enumerate}\setlength{\itemsep}{1.2pt}\setlength{\parskip}{0pt}
  \item \textbf{Packet Delivery Ratio (99.41\%)}: CAPRO attains a $99.41\%$ delivery ratio, significantly outperforming reactive baselines ($91.20\%$) and point prediction ($97.15\%$). By preemptively transferring traffic to stratospheric HAPS prior to convective rain cells, CAPRO maintains link continuity above $98\%$ even during intense precipitation ($100\text{ mm/h}$).
  \item \textbf{Deterministic Tail Latency (6.82 ms)}: CAPRO achieves a 99th-percentile tail latency of $6.82\text{ ms}$, strictly satisfying the sub-10ms URLLC requirement. In contrast, reactive baselines experience large tail delays ($38.45\text{ ms}$) caused by reconnection and packet buffering during sudden outages.
  \item \textbf{Outage Rate Suppression (0.59\%)}: Network outage probability drops from $8.80\%$ to $0.59\%$, representing a $93.3\%$ relative outage reduction across diverse terrestrial clutter and weather states.
  \item \textbf{Handover Stability \& Energy Efficiency}: Risk-adjusted utility suppresses unnecessary handover ping-pong flapping by \textbf{89.3\%} ($1.9$ vs. $18.4$ events/hr), dramatically minimizing signaling overhead and cutting transmission energy per megabit by $56.7\%$ ($2.10$ vs. $4.85\text{ J/Mbit}$).
\end{enumerate}

\section{Research Significance and Applications}
The validated algorithmic contributions provide practical operational benefits across critical 6G deployment domains:
\begin{itemize}\setlength{\itemsep}{1.2pt}\setlength{\parskip}{0pt}
  \item \textbf{Aeronautical \& Maritime Routes}: Guarantees continuous, low-latency broadband for transoceanic flights and deep-sea vessels traversing satellite spot-beam boundaries outside terrestrial radar coverage.
  \item \textbf{Post-Disaster Communications}: Rapidly restores resilient connectivity by dynamically bridging terrestrial first responders to stratospheric HAPS and UAV relays when terrestrial cell towers are incapacitated.
  \item \textbf{Massive Remote IoT}: Provides deterministic telemetry collection for isolated pipeline networks, national power grids, and environmental sensor clusters without local terrestrial base stations.
  \item \textbf{Scalable 6G Architecture}: Demonstrates a modular paradigm separating physical resource allocation (PPO) from cross-tier proactive handover routing (CAPRO) for future mega-constellations.
\end{itemize}

\section{Conclusion and Future Scope}
ORBIT-INTEL 6G presents a comprehensive framework integrating physical-layer deep reinforcement learning with network-layer predictive reliability orchestration for 6G Space-Air-Ground Non-Terrestrial Networks. Combining PPO resource management with distribution-free split-conformal uncertainty calibration enables proactive circumvention of severe atmospheric rain fade and rapid satellite transitions. The 400-trial Monte Carlo evaluation verifies a 99.41\% packet delivery ratio, a 6.82 ms tail latency meeting URLLC requirements, and an 89.3\% reduction in handover ping-pong instability. Future work includes federated reinforcement learning across inter-satellite cross-links, software-defined radio (SDR) hardware-in-the-loop testbed validation, and dynamic spectrum co-existence with terrestrial cellular networks.

\begin{thebibliography}{6}
\fontsize{7.6}{9.2}\selectfont
\setlength{\itemsep}{1.2pt}
\setlength{\parskip}{0pt}
\bibitem{3gpp_tr38811} 3GPP, ``Study on New Radio to support non-terrestrial networks,'' \textit{3GPP TR 38.811}, Rel-15, 2020.
\bibitem{3gpp_tr38821} 3GPP, ``Solutions for NR to support non-terrestrial networks,'' \textit{3GPP TR 38.821}, Rel-16, 2021.
\bibitem{itur_p618} ITU-R, ``Propagation data and prediction methods for Earth-space systems,'' \textit{Rec. ITU-R P.618-13}, 2017.
\bibitem{lin2021jsac} X. Lin et al., ``Multi-beam satellite traffic matching via unconstrained DRL,'' \textit{IEEE JSAC}, 39(4):1120--1134, 2021.
\bibitem{cui2022twc} J. Cui et al., ``Resource allocation in SAGIN using deep reinforcement learning,'' \textit{IEEE TWC}, 21(8):6210--6224, 2022.
\bibitem{angelopoulos2021conformal} A.~N. Angelopoulos and S.~Bates, ``A gentle introduction to conformal prediction,'' \textit{arXiv:2107.07511}, 2021.
\end{thebibliography}

\end{document}
